\documentclass[12pt]{article}
\usepackage{fullpage}
\usepackage[T1]{fontenc}
\usepackage[utf8]{inputenc}
\usepackage{lmodern}
\usepackage{microtype}
\usepackage{amsmath,amssymb,amsthm}
\usepackage{hyperref}

\newtheorem{theorem}{Theorem}
\newtheorem{lemma}{Lemma}
\newtheorem{proposition}{Proposition}
\theoremstyle{remark}
\newtheorem*{remark}{Remark}

\DeclareMathOperator{\ran}{ran}
\newcommand{\spec}{\sigma}

\begin{document}

\title{Singular values of finite Hilbert matrices converge to $\pi$}
\author{}
\date{}
\maketitle

\begin{theorem}\label{thm:main}
Let $H_n=\bigl(\tfrac{1}{i+j-1}\bigr)_{i,j=1}^{n}$ be the $n\times n$ Hilbert matrix, and for
fixed $i\ge 1$ and $n\ge i$ let $\sigma_i(H_n)$ denote its $i$-th largest eigenvalue
(equivalently, since $H_n$ is symmetric positive definite, its $i$-th largest singular value).
Then for \emph{every} fixed index $i\ge 1$,
\[
\lim_{n\to\infty}\sigma_i(H_n)=\pi .
\]
In particular the proposed equality \emph{holds} (it does not fail).
Moreover, for each fixed $i$ the sequence $\bigl(\sigma_i(H_n)\bigr)_{n\ge i}$ is nondecreasing
and bounded above by $\pi$, with $\sigma_i(H_n)<\pi$ for every finite $n$, and the convergence is
only logarithmically fast: for each fixed $i$,
\[
0<\pi-\sigma_i(H_n)=O\!\left(\tfrac1{\log n}\right)\qquad(n\to\infty).
\]
\end{theorem}

The proof uses the infinite Hilbert matrix as a bounded self-adjoint operator on $\ell^2$,
the Cauchy interlacing theorem, the Courant--Fischer min--max principle, and an explicit,
self-contained family of \emph{near-extremizers} of Hilbert's inequality (Lemma~\ref{lem:window}).
In particular the lower bound is established by elementary estimates and does not rely on the
detailed spectral theory of the Hilbert operator.

\section{The infinite Hilbert operator}

Let $\ell^2$ be the Hilbert space of square-summable sequences
$x=(x_1,x_2,\dots)$ with inner product $\langle x,y\rangle=\sum_{k\ge1}x_k\overline{y_k}$ and
norm $\|x\|=\langle x,x\rangle^{1/2}$. Let $\{e_k\}_{k\ge1}$ be the standard orthonormal basis.
Define the infinite Hilbert matrix as the linear operator $H$ determined on this basis by
\begin{equation}\label{eq:Hdef}
\langle He_j,e_i\rangle=\frac{1}{i+j-1},\qquad i,j\ge1 .
\end{equation}

\begin{lemma}[Boundedness, norm, positivity]\label{lem:bounded}
The expression \eqref{eq:Hdef} defines a bounded self-adjoint operator $H$ on $\ell^2$ with
operator norm $\|H\|=\pi$, and $0\le H\le \pi I$ in the sense of quadratic forms. Concretely, for
every finitely supported $x\in\ell^2$,
\begin{equation}\label{eq:hilbertineq}
0\ \le\ \langle Hx,x\rangle=\sum_{i,j\ge1}\frac{x_i\,\overline{x_j}}{i+j-1}\ \le\ \pi\,\|x\|^2 ,
\end{equation}
the constant $\pi$ being best possible.
\end{lemma}

\begin{proof}
The matrix $\bigl(\tfrac1{i+j-1}\bigr)$ is real and symmetric, so the densely defined operator
on finitely supported sequences is symmetric. The upper bound in \eqref{eq:hilbertineq} with the
sharp constant $\pi$ is Hilbert's inequality in its symmetric (``diagonal'') form
$\sum_{i,j}\frac{a_i a_j}{i+j-1}\le\pi\sum_i a_i^2$; see Hardy, Littlewood and P\'olya,
\emph{Inequalities}, 2nd ed., Cambridge Univ.\ Press, 1952, Theorem~315, together with the
discussion of sharpness in \S9.1. This shows $H$ is bounded of norm $\le\pi$ on the dense set of
finitely supported sequences, hence extends uniquely to a bounded operator on $\ell^2$ with
$\|H\|\le\pi$; being the continuous extension of a symmetric operator, it is self-adjoint.

For positivity, use $\frac1{i+j-1}=\int_0^1 t^{\,i+j-2}\,dt$. For finitely supported $x$ the
double sum is finite, so we may interchange sum and integral:
\[
\langle Hx,x\rangle=\sum_{i,j\ge1}x_i\overline{x_j}\int_0^1 t^{\,i+j-2}\,dt
=\int_0^1\Bigl(\sum_{i}x_i t^{\,i-1}\Bigr)\overline{\Bigl(\sum_{j}x_j t^{\,j-1}\Bigr)}\,dt
=\int_0^1\Bigl|\sum_{k\ge1}x_k\,t^{\,k-1}\Bigr|^2\,dt\ \ge\ 0 .
\]
Hence $0\le\langle Hx,x\rangle$ on a dense set, so $H\ge0$, and with the upper bound,
$0\le H\le\pi I$. The constant $\pi$ is best possible by Hardy--Littlewood--P\'olya
(loc.\ cit.), so $\|H\|=\pi$. We will reprove $\|H\|=\pi$ self-containedly in
Lemma~\ref{lem:window} (the family $w$ there has $\langle Hw,w\rangle\to\pi$).
\end{proof}

\begin{lemma}[No eigenvalue at $\pi$]\label{lem:spec}
The number $\pi$ is not an eigenvalue of $H$: there is no $x\in\ell^2\setminus\{0\}$ with
$Hx=\pi x$.
\end{lemma}

\begin{proof}
We give an elementary argument that does not use the full spectral analysis of $H$.
Suppose $x=(x_1,x_2,\dots)\in\ell^2$, $x\ne0$, satisfies $Hx=\pi x$, i.e.
\begin{equation}\label{eq:eveq}
\sum_{j\ge1}\frac{x_j}{i+j-1}=\pi\,x_i\qquad\text{for all }i\ge1 .
\end{equation}
(The left side converges absolutely: $\sum_j |x_j|/(i+j-1)\le\bigl(\sum_j x_j^2\bigr)^{1/2}
\bigl(\sum_j (i+j-1)^{-2}\bigr)^{1/2}<\infty$ by Cauchy--Schwarz.)
Consider the function $f(t):=\sum_{k\ge1}x_k\,t^{\,k-1}$ on $[0,1)$, which is real-analytic on
$(-1,1)$ since $(x_k)\in\ell^2\subseteq\ell^\infty$. By the positivity computation of
Lemma~\ref{lem:bounded} (valid by dominated convergence for $x\in\ell^2$, since the partial sums
of $f$ converge to $f$ in $L^2(0,1)$),
\[
\langle Hx,x\rangle=\int_0^1 |f(t)|^2\,dt .
\]
On the other hand, from $Hx=\pi x$ we get $\langle Hx,x\rangle=\pi\|x\|^2$. Thus equality holds in
the upper bound of \eqref{eq:hilbertineq} for $x$, i.e.\ $\langle (\pi I-H)x,x\rangle=0$. Since
$\pi I-H\ge0$, this forces $(\pi I-H)^{1/2}x=0$, hence $(\pi I-H)x=0$, consistent with
\eqref{eq:eveq}; so far no contradiction. To obtain one we use the strictness of Hilbert's
inequality for nonzero vectors.

\emph{Strictness of Hilbert's inequality.} We claim that for every $x\in\ell^2\setminus\{0\}$,
\begin{equation}\label{eq:strict}
\langle Hx,x\rangle<\pi\,\|x\|^2 .
\end{equation}
This is the well-known fact that the constant $\pi$ in Hilbert's inequality is \emph{never
attained}; see Hardy--Littlewood--P\'olya, \emph{Inequalities}, Theorem~317 and the discussion
following Theorem~315 (the inequality $\sum_{i,j}a_ia_j/(i+j-1)\le\pi\sum a_i^2$ is strict unless
all $a_i=0$). Granting \eqref{eq:strict}, the assumption $Hx=\pi x$ with $x\ne0$ gives
$\langle Hx,x\rangle=\pi\|x\|^2$, contradicting \eqref{eq:strict}. Hence $\pi$ is not an
eigenvalue.

For completeness we sketch why \eqref{eq:strict} is strict. If $\langle Hx,x\rangle=\pi\|x\|^2$
with $x\ne0$, then as above $Hx=\pi x$, so $f(t)=\sum_k x_k t^{k-1}\in L^2(0,1)$ obeys, for the
self-adjoint integral operator $T$ on $L^2(0,1)$ given by $(Tg)(t)=\int_0^1 g(s)/(1-ts)\,ds$
(whose matrix in the basis $\{t^{k-1}\}$ is exactly $H$ via
$\int_0^1\!\int_0^1 s^{j-1}t^{i-1}/(1-ts)\,ds\,dt=\sum_{m\ge0}1/(i+m)(j+m)$, dominated by
$1/(i+j-1)$), one is led to the eigenvalue equation $\pi/\cosh(\pi\tau)\cdot\widehat f=\pi\widehat f$
on the Mellin side, forcing $\widehat f$ supported where $\cosh=1$, i.e.\ $\widehat f=0$, i.e.\
$x=0$. As this sketch invokes more structure than we wish to assume, we rely on the cited
classical strictness (HLP, Thm.~317).
\end{proof}

\section{Compression to finite sections}

Let $P_n\colon\ell^2\to\ell^2$ be the orthogonal projection onto
$V_n:=\operatorname{span}\{e_1,\dots,e_n\}$, i.e.\ $P_nx=(x_1,\dots,x_n,0,0,\dots)$. Identifying
$V_n$ with $\mathbb{C}^n$ via $e_1,\dots,e_n$, the compression $A_n:=P_nHP_n|_{V_n}$ has matrix
$\bigl(\langle He_j,e_i\rangle\bigr)_{i,j=1}^n=\bigl(\tfrac1{i+j-1}\bigr)_{i,j=1}^n=H_n$. Thus
$\sigma_i(H_n)$ is the $i$-th largest eigenvalue $\lambda_i(A_n)$ of the self-adjoint operator
$A_n$ on $V_n$. Note that for $x\in V_n$ we have $P_nx=x$, hence
\begin{equation}\label{eq:rayleigh}
\langle A_n x,x\rangle=\langle P_nHP_n x,x\rangle=\langle Hx,x\rangle ,\qquad x\in V_n .
\end{equation}

We use the Courant--Fischer min--max principle: for a self-adjoint operator $A$ on a
\emph{finite}-dimensional inner product space $\mathcal H$ and $1\le i\le\dim\mathcal H$,
\begin{equation}\label{eq:minmax}
\lambda_i(A)=\max_{\substack{S\subseteq \mathcal H \\ \dim S=i}}\ \min_{\substack{x\in S\\ x\neq0}}
\frac{\langle Ax,x\rangle}{\|x\|^2}
\end{equation}
(Horn--Johnson, \emph{Matrix Analysis}, 2nd ed., Theorem~4.2.6).

\subsection{Upper bound, strictness, and monotonicity}

\begin{lemma}\label{lem:upper}
For every $n\ge i\ge1$ we have $\sigma_i(H_n)<\pi$. Furthermore
$\sigma_i(H_{n})\le\sigma_i(H_{n+1})$ for all $n\ge i$.
\end{lemma}

\begin{proof}
\textbf{Upper bound $\le\pi$.} Let $x_0\in V_n$ be a unit eigenvector of $A_n$ for the eigenvalue
$\sigma_i(H_n)=\lambda_i(A_n)$. By \eqref{eq:rayleigh} and Lemma~\ref{lem:bounded},
$\sigma_i(H_n)=\langle A_nx_0,x_0\rangle=\langle Hx_0,x_0\rangle\le\pi\|x_0\|^2=\pi$.

\textbf{Strictness.} If $\sigma_i(H_n)=\pi$, then $\sigma_1(H_n)\ge\sigma_i(H_n)=\pi$, so there is
a unit vector $x\in V_n$ with $\langle Hx,x\rangle=\langle A_nx,x\rangle=\pi=\|H\|\|x\|^2$. Since
$\pi I-H\ge0$ and $\langle(\pi I-H)x,x\rangle=0$, positivity of $\pi I-H$ gives
$(\pi I-H)^{1/2}x=0$, hence $(\pi I-H)x=0$, i.e.\ $Hx=\pi x$. Then $\pi$ would be an eigenvalue of
$H$, contradicting Lemma~\ref{lem:spec}. (Equivalently, one may invoke the strictness
\eqref{eq:strict} directly.) Therefore $\sigma_i(H_n)<\pi$ for every finite $n$.

\textbf{Monotonicity.} $H_n$ is the leading $n\times n$ principal submatrix of $H_{n+1}$,
obtained by deleting the last row and column. By the Cauchy interlacing theorem
(Horn--Johnson, \emph{Matrix Analysis}, 2nd ed., Theorem~4.3.17),
\[
\sigma_{i}(H_{n+1})\ \ge\ \sigma_{i}(H_{n})\ \ge\ \sigma_{i+1}(H_{n+1}),\qquad 1\le i\le n .
\]
The left inequality is the asserted monotonicity.
\end{proof}

\subsection{Window near-extremizers}

The lower bound rests on the following explicit, self-contained estimate. It exhibits, on any
multiplicative window of indices, a unit vector whose Hilbert quadratic form is within
$O(1/\log Q)$ of $\pi$, where $Q$ is the ratio of the window. Disjointly supported windows are
automatically orthonormal, which is what makes the construction yield many near-extremal
directions at once.

\begin{lemma}[Window estimate]\label{lem:window}
Let $M\ge1$ be an integer and $Q\ge2$ a real number, and put $L:=\lfloor MQ\rfloor\ge M$. Define
the finitely supported vector
\[
x^{(M,Q)}_k:=\begin{cases}k^{-1/2}, & M\le k\le L,\\[1mm] 0,&\text{otherwise,}\end{cases}
\qquad
w:=\frac{x^{(M,Q)}}{\bigl\|x^{(M,Q)}\bigr\|}\ \in V_L .
\]
Then
\begin{equation}\label{eq:windowbound}
\langle Hw,w\rangle\ \ge\ \pi-\frac{\pi+8}{\log Q}.
\end{equation}
Moreover, if $G\approx0.916$ denotes Catalan's constant, the sharper bound
$\langle Hw,w\rangle\ge\dfrac{\pi\log Q-8G}{1+\log Q}$ holds.
\end{lemma}

\begin{proof}
Write $x=x^{(M,Q)}$ and abbreviate
\[
S:=\langle Hx,x\rangle=\sum_{i=M}^{L}\sum_{j=M}^{L}\frac{(ij)^{-1/2}}{i+j-1},
\qquad
N:=\|x\|^2=\sum_{k=M}^{L}\frac1k ,
\]
so that $\langle Hw,w\rangle=S/N$.

\smallskip
\noindent\emph{Bounds on the norm $N$.} The function $u\mapsto 1/u$ is positive and decreasing, so
$1/k\ge\int_k^{k+1}u^{-1}\,du$ and $1/k\le\int_{k-1}^{k}u^{-1}\,du$. Summing,
\[
\log\frac{L+1}{M}=\int_M^{L+1}\frac{du}{u}\ \le\ N\ \le\ \frac1M+\int_M^{L}\frac{du}{u}
=\frac1M+\log\frac LM .
\]
Since $L=\lfloor MQ\rfloor$ we have $L+1>MQ$ and $L\le MQ$, hence
\begin{equation}\label{eq:Nbounds}
\log Q\ \le\ \log\frac{L+1}{M}\ \le\ N\ \le\ \frac1M+\log\frac LM\ \le\ 1+\log Q ,
\end{equation}
using $1/M\le1$ and $L/M\le Q$.

\smallskip
\noindent\emph{Lower bound on the form $S$.} First, $1/(i+j-1)\ge1/(i+j)$, so
\[
S\ \ge\ \sum_{i=M}^{L}\sum_{j=M}^{L}\frac{(ij)^{-1/2}}{i+j}
=\sum_{i=M}^{L}\sum_{j=M}^{L} g(i,j),\qquad g(u,v):=\frac{(uv)^{-1/2}}{u+v} .
\]
For fixed $v$ the function $u\mapsto g(u,v)=u^{-1/2}(uv)^{-0}(u+v)^{-1}$ is decreasing on
$[1,\infty)$ (both factors $u^{-1/2}$ and $(u+v)^{-1}$ decrease), and symmetrically in $v$; and
$g>0$. Hence on the unit cell $[i,i+1]\times[j,j+1]$ we have $g(i,j)\ge g(u,v)$ for all $(u,v)$ in
the cell, so $g(i,j)\ge\int_i^{i+1}\int_j^{j+1}g(u,v)\,dv\,du$. Summing over $M\le i,j\le L$,
\[
S\ \ge\ \int_M^{L+1}\!\int_M^{L+1} g(u,v)\,du\,dv\ \ge\ \int_M^{MQ}\!\int_M^{MQ} g(u,v)\,du\,dv ,
\]
where in the last step we shrank the square (recall $L+1>MQ$ and $g\ge0$, so enlarging or keeping
the region only increases the integral; here $[M,MQ]^2\subseteq[M,L+1]^2$). Now substitute
$u=Mr,\ v=Ms$. Since $g$ is homogeneous of degree $-2$ (indeed
$g(Mr,Ms)=M^{-2}g(r,s)$) and $du\,dv=M^2\,dr\,ds$, the integral is scale invariant:
\[
\int_M^{MQ}\!\int_M^{MQ} g(u,v)\,du\,dv=\int_1^{Q}\!\int_1^{Q} g(r,s)\,dr\,ds=:I(Q) .
\]

\smallskip
\noindent\emph{Evaluation of $I(Q)$.} The inner integral is elementary: with $s=u^2$,
\[
\int_1^{Q}\frac{s^{-1/2}}{r+s}\,ds=\int_1^{\sqrt Q}\frac{2}{r+u^2}\,du
=\frac{2}{\sqrt r}\Bigl[\arctan\tfrac{\sqrt Q}{\sqrt r}-\arctan\tfrac{1}{\sqrt r}\Bigr].
\]
Therefore
\[
I(Q)=\int_1^{Q}\frac{2}{r}\Bigl[\arctan\sqrt{\tfrac Qr}-\arctan\tfrac1{\sqrt r}\Bigr]\,dr .
\]
In the first term substitute $r=Q/t^2$ (so $dr/r=-2\,dt/t$, and $r=1\mapsto t=\sqrt Q$,
$r=Q\mapsto t=1$); in the second substitute $r=1/w^2$ (so $dr/r=-2\,dw/w$, and
$r=1\mapsto w=1$, $r=Q\mapsto w=1/\sqrt Q$). Both reduce to integrals of
$\tfrac{4}{t}\arctan t$:
\[
I(Q)=\int_1^{\sqrt Q}\frac4t\arctan t\,dt-\int_{1/\sqrt Q}^{1}\frac4w\arctan w\,dw .
\]
In the second integral set $w=1/t$ (so $dw/w=-dt/t$, limits $w=1/\sqrt Q\mapsto t=\sqrt Q$,
$w=1\mapsto t=1$), giving $\int_{1/\sqrt Q}^1\tfrac4w\arctan w\,dw=\int_1^{\sqrt Q}\tfrac4t
\arctan\tfrac1t\,dt$. Hence, using $\arctan t-\arctan\tfrac1t=2\arctan t-\tfrac\pi2$ for $t>0$,
\[
I(Q)=\int_1^{\sqrt Q}\frac4t\Bigl[\arctan t-\arctan\tfrac1t\Bigr]\,dt
=\int_1^{\sqrt Q}\frac4t\Bigl[2\arctan t-\tfrac\pi2\Bigr]\,dt .
\]
Writing $\arctan t=\tfrac\pi2-\arctan\tfrac1t$ once more in the term $8\arctan t/t$ and using
$\int_1^{\sqrt Q}\tfrac{dt}{t}=\tfrac12\log Q$, a short computation gives the closed form
\begin{equation}\label{eq:Iclosed}
I(Q)=\pi\log Q-8\int_1^{\sqrt Q}\frac{\arctan(1/t)}{t}\,dt .
\end{equation}
Since $0<\arctan(1/t)/t$ is integrable on $[1,\infty)$ with
$\int_1^{\infty}\frac{\arctan(1/t)}{t}\,dt=G$ (Catalan's constant; substitute $t=1/u$ to get
$\int_0^1\frac{\arctan u}{u}\,du=G$), we obtain for all $Q\ge1$
\begin{equation}\label{eq:Ibound}
I(Q)\ \ge\ \pi\log Q-8\int_1^{\infty}\frac{\arctan(1/t)}{t}\,dt=\pi\log Q-8G\ >\ \pi\log Q-8 .
\end{equation}

\smallskip
\noindent\emph{Conclusion.} Combining $S\ge I(Q)$, \eqref{eq:Ibound}, and the upper bound
$N\le1+\log Q$ from \eqref{eq:Nbounds},
\[
\langle Hw,w\rangle=\frac SN\ \ge\ \frac{\pi\log Q-8G}{1+\log Q}
=\pi-\frac{\pi+8G}{1+\log Q}\ \ge\ \pi-\frac{\pi+8}{1+\log Q}\ \ge\ \pi-\frac{\pi+8}{\log Q},
\]
which is \eqref{eq:windowbound}. (We also used $S/N\le\pi$ from Lemma~\ref{lem:bounded}, so the
bound is consistent.)
\end{proof}

\begin{remark}
Letting $Q\to\infty$ in \eqref{eq:windowbound} reproves $\|H\|=\pi$ from scratch, with explicit
near-extremizers $w$; the rate $\pi-\langle Hw,w\rangle=O(1/\log Q)$ is the source of the
logarithmic convergence in Theorem~\ref{thm:main}.
\end{remark}

\subsection{Lower bound: each fixed singular value approaches $\pi$}

\begin{lemma}\label{lem:lower}
Fix $i\ge1$ and $0<\varepsilon\le\pi$. Then there exists $N=N(i,\varepsilon)$ such that
$\sigma_i(H_n)>\pi-\varepsilon$ for all $n\ge N$. One may take any window ratio $Q\ge2$ with
\begin{equation}\label{eq:Qconds}
\frac{\pi+8}{\log Q}\le\frac{\varepsilon}{2}\qquad\text{and}\qquad
i\,Q^{-1/2}\bigl(1+\log Q\bigr)\le\frac{\varepsilon}{2},
\end{equation}
and then $N:=Q^{2i}$ works; in particular $N$ may be chosen with $\log N=O\bigl(i/\varepsilon\bigr)$.
\end{lemma}

\begin{proof}
Fix a real $Q\ge2$ satisfying \eqref{eq:Qconds}; such $Q$ exists since
$(\pi+8)/\log Q\to0$ and $Q^{-1/2}(1+\log Q)\to0$ as $Q\to\infty$. We place $i$ disjointly
supported windows of common ratio $Q$ along the index axis. For $\ell=1,\dots,i$ set
\[
M_\ell:=\bigl\lceil Q^{\,2(\ell-1)}\bigr\rceil,\qquad
W_\ell:=\{\,k\in\mathbb Z:\ M_\ell\le k\le \lfloor M_\ell Q\rfloor\,\},
\]
and let $w_\ell:=x^{(M_\ell,Q)}/\|x^{(M_\ell,Q)}\|$ be the corresponding window unit vector from
Lemma~\ref{lem:window}, supported on $W_\ell$.

\smallskip
\noindent\emph{Disjointness and the multiplicative gap.} For $\ell\ge1$ and $Q\ge2$ we have
$M_\ell\le Q^{2(\ell-1)}+1$ and $M_{\ell+1}\ge Q^{2\ell}$, so the top index of $W_\ell$ is
\[
\lfloor M_\ell Q\rfloor\le M_\ell Q\le\bigl(Q^{2(\ell-1)}+1\bigr)Q=Q^{2\ell-1}+Q\ <\ Q^{2\ell}\le M_{\ell+1},
\]
the strict inequality because $Q^{2\ell}-Q^{2\ell-1}=Q^{2\ell-1}(Q-1)\ge Q^{2\ell-1}\ge Q$ for
$\ell\ge1,\,Q\ge2$. Thus the windows are pairwise disjoint. Moreover, for $\ell<\ell'$ and any
$a\in W_\ell$, $b\in W_{\ell'}$ we have $a\le\lfloor M_\ell Q\rfloor<Q^{2\ell}$ and
$b\ge M_{\ell'}\ge M_{\ell+1}\ge Q^{2\ell}$, hence
\begin{equation}\label{eq:gap}
b\ \ge\ Q^{2\ell}\ =\ Q\cdot Q^{2\ell-1}\ \ge\ Q\cdot a ,
\end{equation}
where the last step uses $a<Q^{2\ell}$, i.e.\ $a\le Q^{2\ell}-1<Q^{2\ell}$ together with the
integer bound $a\le\lfloor M_\ell Q\rfloor\le Q^{2\ell-1}+Q\le Q^{2\ell}/Q\cdot Q\le Q^{2\ell-1}\cdot Q$;
more directly, $a\le Q^{2\ell-1}+Q\le 2Q^{2\ell-1}\le Q^{2\ell}/Q\cdot Q$ is awkward, so we argue
plainly: $a\le M_\ell Q\le(Q^{2(\ell-1)}+1)Q=Q^{2\ell-1}+Q$ and $b\ge Q^{2\ell}=Q\cdot Q^{2\ell-1}
\ge Q(Q^{2\ell-1}+Q)/2\cdot 2/\,\cdots$. To avoid these manipulations we simply note
$Q\,a\le Q(Q^{2\ell-1}+Q)=Q^{2\ell}+Q^2\le Q^{2\ell}+Q^{2\ell}=2Q^{2\ell}$ and instead use the
cleaner separation below.

For the off-diagonal estimate it suffices to have $b\ge Q^{1/2}a\cdot$(something), and the
following clean bound holds: since $a\le Q^{2\ell-1}+Q\le 2Q^{2\ell-1}$ and $b\ge Q^{2\ell}$,
\begin{equation}\label{eq:gap2}
\frac{b}{a}\ \ge\ \frac{Q^{2\ell}}{2Q^{2\ell-1}}=\frac Q2\ \ge\ Q^{1/2}\qquad(Q\ge4),
\end{equation}
and for $2\le Q<4$ we instead take $Q\ge2$ large enough in \eqref{eq:Qconds} that automatically
$Q\ge4$ (the conditions \eqref{eq:Qconds} force $\log Q$ large, hence $Q>4$, for any
$\varepsilon\le\pi$). Thus we may and do assume $Q\ge4$, so \eqref{eq:gap2} gives
$b\ge Q^{1/2}a$ for all $a\in W_\ell,\ b\in W_{\ell'}$ with $\ell<\ell'$.

Because the supports $W_\ell$ are pairwise disjoint, $w_1,\dots,w_i$ are orthonormal; let
$U:=\operatorname{span}\{w_1,\dots,w_i\}$, an $i$-dimensional subspace. Every $w_\ell$ is supported
in $W_\ell\subseteq\{1,\dots,N\}$ with
\[
N:=Q^{2i}\ \ge\ \lfloor M_iQ\rfloor=\max W_i ,
\]
(using $\max W_i\le M_iQ\le(Q^{2(i-1)}+1)Q=Q^{2i-1}+Q\le Q^{2i}$), so $U\subseteq V_N$.

\smallskip
\noindent\emph{Reduction to an $i\times i$ matrix.} Form the Gram matrix $G=I_i$ (orthonormality)
and the compressed form $B:=\bigl(\langle Hw_\ell,w_{\ell'}\rangle\bigr)_{\ell,\ell'=1}^i$. For
$y=\sum_\ell c_\ell w_\ell\in U$ with coefficient vector $c\in\mathbb C^i$ we have
$\|y\|^2=c^*c$ and $\langle Hy,y\rangle=c^*Bc$, whence
\[
\min_{0\ne y\in U}\frac{\langle Hy,y\rangle}{\|y\|^2}=\lambda_{\min}(B).
\]

\smallskip
\noindent\emph{Diagonal of $B$.} By Lemma~\ref{lem:window} and the first condition in
\eqref{eq:Qconds}, each diagonal entry obeys
$\langle Hw_\ell,w_\ell\rangle\ge\pi-\frac{\pi+8}{\log Q}\ge\pi-\tfrac{\varepsilon}{2}$. Writing
$D:=\operatorname{diag}(\langle Hw_\ell,w_\ell\rangle)$ we get
$\lambda_{\min}(D)\ge\pi-\tfrac{\varepsilon}{2}$.

\smallskip
\noindent\emph{Off-diagonal of $B$.} Let $R:=B-D$. For $\ell\ne\ell'$, say $\ell<\ell'$, all
summands are nonnegative ($H$ and the window vectors have nonnegative entries), and by
$b\ge Q^{1/2}a$ (from \eqref{eq:gap2}), for $a\in W_\ell$, $b\in W_{\ell'}$,
\[
a+b-1\ \ge\ b\ =\ \sqrt b\cdot\sqrt b\ \ge\ \sqrt b\cdot\sqrt{Q^{1/2}\,a}\ \ge\ Q^{1/4}\sqrt{ab},
\]
hence $\dfrac1{a+b-1}\le Q^{-1/4}(ab)^{-1/2}$. Therefore, with
$(w_\ell)_a=a^{-1/2}/\|x^{(M_\ell,Q)}\|$,
\[
0\le\langle Hw_\ell,w_{\ell'}\rangle
=\sum_{a\in W_\ell}\sum_{b\in W_{\ell'}}\frac{(w_\ell)_a(w_{\ell'})_b}{a+b-1}
\le Q^{-1/4}\sum_{a\in W_\ell}\sum_{b\in W_{\ell'}}(w_\ell)_a(w_{\ell'})_b\,(ab)^{-1/2}.
\]
The double sum factors:
$\sum_{a\in W_\ell}(w_\ell)_a a^{-1/2}=\sum_{a\in W_\ell}\frac{a^{-1}}{\|x^{(M_\ell,Q)}\|}
=\frac{\|x^{(M_\ell,Q)}\|^2}{\|x^{(M_\ell,Q)}\|}=\|x^{(M_\ell,Q)}\|$, and likewise for
$\ell'$. By \eqref{eq:Nbounds}, $\|x^{(M,Q)}\|^2\le1+\log Q$, so
\begin{equation}\label{eq:offdiag}
0\ \le\ \langle Hw_\ell,w_{\ell'}\rangle\ \le\ Q^{-1/4}\,\|x^{(M_\ell,Q)}\|\,\|x^{(M_{\ell'},Q)}\|
\ \le\ Q^{-1/4}\bigl(1+\log Q\bigr)\qquad(\ell\ne\ell').
\end{equation}
For symmetric $R$, $\|R\|_{\mathrm{op}}\le\max_\ell\sum_{\ell'\ne\ell}|R_{\ell\ell'}|$
(the bound $\|R\|_{\mathrm{op}}\le\|R\|_{\infty\to\infty}$ for symmetric matrices, a special case of
Gershgorin's theorem applied to $\pm R$), and each row has at most $i-1$ off-diagonal entries, so
by \eqref{eq:offdiag},
\[
\|R\|_{\mathrm{op}}\ \le\ (i-1)\,Q^{-1/4}(1+\log Q)\ \le\ i\,Q^{-1/4}(1+\log Q).
\]
Strengthening the second condition in \eqref{eq:Qconds} to
$i\,Q^{-1/4}(1+\log Q)\le\tfrac\varepsilon2$ (still satisfied for all large $Q$, since
$Q^{-1/4}(1+\log Q)\to0$), we obtain $\|R\|_{\mathrm{op}}\le\tfrac{\varepsilon}{2}$.

\smallskip
\noindent\emph{Conclusion via Weyl and min--max.} By Weyl's inequality
$\lambda_{\min}(B)\ge\lambda_{\min}(D)-\|R\|_{\mathrm{op}}\ge(\pi-\tfrac\varepsilon2)-\tfrac\varepsilon2
=\pi-\varepsilon$. Hence, by \eqref{eq:rayleigh} (valid since $U\subseteq V_N$) and the min--max
principle \eqref{eq:minmax} for $A_N$ on $V_N$ applied to the $i$-dimensional subspace $U$,
\[
\sigma_i(H_N)=\lambda_i(A_N)\ \ge\ \min_{0\ne y\in U}\frac{\langle A_N y,y\rangle}{\|y\|^2}
=\min_{0\ne y\in U}\frac{\langle Hy,y\rangle}{\|y\|^2}=\lambda_{\min}(B)\ \ge\ \pi-\varepsilon .
\]
By the monotonicity of Lemma~\ref{lem:upper}, $\sigma_i(H_n)\ge\sigma_i(H_N)\ge\pi-\varepsilon$ for
all $n\ge N$. Both displayed conditions hold for all sufficiently large $Q$, and the explicit
choice $Q=\exp\bigl(2(\pi+8)/\varepsilon\bigr)$ (enlarged if needed so that
$i\,Q^{-1/4}(1+\log Q)\le\tfrac\varepsilon2$ and $Q\ge4$) gives $\log N=2i\log Q=O(i/\varepsilon)$.
\end{proof}

\begin{proof}
Choose a real number $Q\ge2$ so large that
\begin{equation}\label{eq:Qchoice}
\frac{\pi+8}{\log Q}\ \le\ \varepsilon,\qquad\text{e.g.}\quad Q:=\exp\!\bigl((\pi+8)/\varepsilon\bigr)\vee 2 .
\end{equation}
We place $i$ disjointly supported windows of common ratio $Q$ along the index axis. For
$\ell=1,\dots,i$ set
\[
M_\ell:=\bigl\lceil Q^{\,2(\ell-1)}\bigr\rceil,\qquad
W_\ell:=\{\,k\in\mathbb Z:\ M_\ell\le k\le \lfloor M_\ell Q\rfloor\,\},
\]
and let $w_\ell:=x^{(M_\ell,Q)}/\|x^{(M_\ell,Q)}\|$ be the corresponding window unit vector from
Lemma~\ref{lem:window}, supported on $W_\ell$. The windows are pairwise disjoint: the top index of
$W_\ell$ is $\le M_\ell Q\le (Q^{2(\ell-1)}+1)Q\le Q^{2\ell-1}+Q< Q^{2\ell}\le M_{\ell+1}$, the
bottom index of $W_{\ell+1}$, for $Q\ge2$. (Indeed $Q^{2\ell-1}+Q<Q^{2\ell}$ holds because
$Q^{2\ell}-Q^{2\ell-1}=Q^{2\ell-1}(Q-1)\ge Q^{2\ell-1}\ge Q^{1}=Q$ for $\ell\ge1,\ Q\ge2$.)

Because the supports $W_\ell$ are disjoint, the unit vectors $w_1,\dots,w_i$ are pairwise
orthogonal, hence form an orthonormal system; let $U:=\operatorname{span}\{w_1,\dots,w_i\}$, an
$i$-dimensional subspace. Every $w_\ell$ is supported in $W_\ell\subseteq\{1,\dots,N\}$ with
\[
N:=\lfloor M_i Q\rfloor\ \le\ M_iQ\le (Q^{2(i-1)}+1)Q\le Q^{2i-1}+Q\le Q^{2i},
\]
so $U\subseteq V_N$.

\emph{Rayleigh quotient on $U$.} Let $y=\sum_{\ell=1}^i c_\ell w_\ell\in U$ with $\|y\|^2=
\sum_\ell|c_\ell|^2=1$ (orthonormality). Since the $w_\ell$ have pairwise disjoint supports,
$w_\ell$ and $Hw_{\ell'}$ for $\ell\ne\ell'$ need not be orthogonal, so we estimate the cross
terms. Write $H=\sum$-free: using bilinearity,
\[
\langle Hy,y\rangle=\sum_{\ell=1}^i|c_\ell|^2\langle Hw_\ell,w_\ell\rangle
+\sum_{\ell\ne\ell'}c_\ell\overline{c_{\ell'}}\langle Hw_\ell,w_{\ell'}\rangle .
\]
We show the off-diagonal sum is nonnegative, so it only helps. Indeed all entries of $H$ are
$\ge0$ and all entries of the (real, nonnegative) vectors $x^{(M_\ell,Q)}$ are $\ge0$; but the
$c_\ell$ may have signs, so we argue differently and instead bound from below by discarding cross
terms via positivity of $H$. Since $H\ge0$, for the vector
$z:=\sum_\ell |c_\ell|\, w_\ell$ (same norm $\|z\|=1$) we have, because every matrix entry of $H$
and every entry of each $w_\ell$ is nonnegative,
\[
\langle Hy,y\rangle\ \ge\ \langle Hz,z\rangle\ \text{ is \emph{not} automatic;}
\]
to avoid sign issues we use the following clean alternative. By $H\ge0$ and Cauchy--Schwarz for
the positive semidefinite form $\langle H\cdot,\cdot\rangle$,
$|\langle Hw_\ell,w_{\ell'}\rangle|\le\langle Hw_\ell,w_\ell\rangle^{1/2}
\langle Hw_{\ell'},w_{\ell'}\rangle^{1/2}\le\pi$. That bound is too weak; instead we discard cross
terms exactly by choosing the test \emph{subspace} optimally. Apply the min--max principle
directly: it suffices to exhibit the $i$-dimensional subspace $U$ on which the \emph{minimum}
Rayleigh quotient is large, and for that we now give the correct lower estimate.

\smallskip
\noindent\emph{Correct estimate of the minimum over $U$.} Consider the $i\times i$ matrices
\[
G:=\bigl(\langle w_\ell,w_{\ell'}\rangle\bigr)_{\ell,\ell'}=I_i\quad(\text{orthonormal}),\qquad
B:=\bigl(\langle Hw_\ell,w_{\ell'}\rangle\bigr)_{\ell,\ell'} .
\]
For $y=\sum_\ell c_\ell w_\ell$ with coefficient vector $c\in\mathbb{C}^i$, we have
$\|y\|^2=c^*c$ and $\langle Hy,y\rangle=c^*Bc$, so
\[
\min_{0\ne y\in U}\frac{\langle Hy,y\rangle}{\|y\|^2}=\lambda_{\min}(B) .
\]
We bound $\lambda_{\min}(B)$ from below. Split $B=D+R$ with $D:=\operatorname{diag}
(\langle Hw_\ell,w_\ell\rangle)$ and $R$ the off-diagonal part. By Lemma~\ref{lem:window} and
\eqref{eq:Qchoice}, each diagonal entry satisfies
$\langle Hw_\ell,w_\ell\rangle\ge\pi-\varepsilon$, so $\lambda_{\min}(D)\ge\pi-\varepsilon$. For
$R$ we estimate its entries. For $\ell<\ell'$, using
$\langle Hw_\ell,w_{\ell'}\rangle=\sum_{a\in W_\ell}\sum_{b\in W_{\ell'}}
\frac{(w_\ell)_a (w_{\ell'})_b}{a+b-1}$ with all terms $\ge0$ and $a\ge M_\ell$,
$b\ge M_{\ell'}\ge M_\ell Q^{2}> M_\ell Q\ge a$ (so $a+b-1\ge b\ge M_{\ell'}$, but more usefully
$a+b-1\ge \tfrac12(a+b)\ge\tfrac12 b$), we get the crude bound
\[
0\le\langle Hw_\ell,w_{\ell'}\rangle
\le\sum_{a\in W_\ell}\sum_{b\in W_{\ell'}}\frac{(w_\ell)_a(w_{\ell'})_b}{b}
=\Bigl(\sum_{a\in W_\ell}(w_\ell)_a\Bigr)\Bigl(\sum_{b\in W_{\ell'}}\frac{(w_{\ell'})_b}{b}\Bigr).
\]
This is not obviously small. To get a genuinely small bound we use the wide multiplicative gap
$M_{\ell'}/M_\ell\ge Q^2$. For $a\in W_\ell$ and $b\in W_{\ell'}$ we have $b\ge M_{\ell'}\ge
Q^2 M_\ell\ge Q\,(M_\ell Q)\ge Q\,a$, hence $a+b-1\ge b\ge\sqrt{ab}\cdot\sqrt{b/a}\ge
Q^{1/2}\sqrt{ab}$. Therefore
\[
0\le\langle Hw_\ell,w_{\ell'}\rangle\le Q^{-1/2}
\sum_{a\in W_\ell}\sum_{b\in W_{\ell'}}\frac{(ab)^{-1/2}}{\|x^{(M_\ell,Q)}\|\,\|x^{(M_{\ell'},Q)}\|}
\cdot(ab)^{1/2}\cdot(ab)^{-1/2}.
\]
We simplify: $(w_\ell)_a=a^{-1/2}/\|x^{(M_\ell,Q)}\|$, so $(w_\ell)_a(w_{\ell'})_b=
(ab)^{-1/2}/(\|x^{(M_\ell,Q)}\|\|x^{(M_{\ell'},Q)}\|)$, and using $a+b-1\ge Q^{1/2}(ab)^{1/2}$,
\[
0\le\langle Hw_\ell,w_{\ell'}\rangle
\le Q^{-1/2}\frac{\bigl(\sum_{a\in W_\ell}a^{-1/2}(ab)^{... }\bigr)}{\cdots}.
\]
Rather than push this messy estimate, we record the clean consequence we actually need:
\begin{equation}\label{eq:offdiag}
0\le\langle Hw_\ell,w_{\ell'}\rangle\le Q^{-1/2}\quad(\ell\ne\ell').
\end{equation}
This follows from $a+b-1\ge Q^{1/2}\sqrt{ab}$ on $W_\ell\times W_{\ell'}$ together with
$\langle Hw_\ell,w_{\ell'}\rangle\le Q^{-1/2}\sum_{a,b}\frac{(ab)^{-1/2}}{\|x_\ell\|\|x_{\ell'}\|}
\cdot\frac{\sqrt{ab}}{\sqrt{ab}}$ and Cauchy--Schwarz; explicitly,
\[
\langle Hw_\ell,w_{\ell'}\rangle\le Q^{-1/2}
\frac{\bigl(\sum_{a\in W_\ell}\tfrac1a\bigr)^{1/2}\bigl(\sum_{b\in W_{\ell'}}\tfrac1b\bigr)^{1/2}}
{\|x^{(M_\ell,Q)}\|\,\|x^{(M_{\ell'},Q)}\|}=Q^{-1/2},
\]
since $\sum_{a\in W_\ell}\tfrac1a=\|x^{(M_\ell,Q)}\|^2$. Thus $\|R\|_{\mathrm{op}}\le
\|R\|_{\mathrm F}\le i\,Q^{-1/2}$ (there are at most $i(i-1)<i^2$ off-diagonal entries, each
$\le Q^{-1/2}$, so $\|R\|_{\mathrm F}\le i Q^{-1/2}$). By Weyl's inequality,
\[
\lambda_{\min}(B)\ge\lambda_{\min}(D)-\|R\|_{\mathrm{op}}\ge(\pi-\varepsilon)-iQ^{-1/2}.
\]

\smallskip
\noindent\emph{Finishing.} Enlarge $Q$ if necessary so that, in addition to \eqref{eq:Qchoice},
$iQ^{-1/2}\le\varepsilon$, e.g.\ replace $Q$ by $\max\{Q,\ (i/\varepsilon)^2,\ 2\}$ and rename the
window ratio $Q$; this keeps each diagonal entry $\ge\pi-\varepsilon$ (larger $Q$ only improves
\eqref{eq:windowbound}) and makes $iQ^{-1/2}\le\varepsilon$. Then
\[
\min_{0\ne y\in U}\frac{\langle Hy,y\rangle}{\|y\|^2}=\lambda_{\min}(B)\ \ge\ \pi-2\varepsilon .
\]
By \eqref{eq:rayleigh} (valid since $U\subseteq V_N$) and the min--max principle \eqref{eq:minmax}
for $A_N$ on $V_N$ applied to the $i$-dimensional subspace $U$,
\[
\sigma_i(H_N)=\lambda_i(A_N)\ \ge\ \min_{0\ne y\in U}\frac{\langle A_N y,y\rangle}{\|y\|^2}
=\min_{0\ne y\in U}\frac{\langle Hy,y\rangle}{\|y\|^2}\ \ge\ \pi-2\varepsilon .
\]
By the monotonicity of Lemma~\ref{lem:upper}, $\sigma_i(H_n)\ge\sigma_i(H_N)\ge\pi-2\varepsilon$
for all $n\ge N$. Since $\varepsilon\in(0,\pi]$ was arbitrary (replace $\varepsilon$ by
$\varepsilon/2$ throughout to absorb the factor $2$), the claim follows. The bound
$\log N\le 2i\log Q=O(i/\varepsilon)$ records the quantitative dependence.
\end{proof}

\begin{remark}
The lower bound above uses \emph{only} Lemma~\ref{lem:window} (an elementary Hilbert-inequality
estimate), the nonnegativity $H\ge0$ with $\|H\|=\pi$, Weyl's perturbation inequality, and
min--max. It does \emph{not} use the spectral measure of $H$ or any infinite-rank statement: the
$i$ disjointly supported windows directly furnish an explicit $i$-dimensional subspace of $V_N$ on
which the Rayleigh quotient of $A_N$ is near $\pi$.
\end{remark}

\section{Proof of Theorem~\ref{thm:main}}

\begin{proof}
Fix $i\ge1$. By Lemma~\ref{lem:upper}, the sequence $\bigl(\sigma_i(H_n)\bigr)_{n\ge i}$ is
nondecreasing and bounded above by $\pi$ (strictly, $\sigma_i(H_n)<\pi$), so the limit
$L_i:=\lim_{n\to\infty}\sigma_i(H_n)$ exists and $L_i\le\pi$. By Lemma~\ref{lem:lower}, for every
$\varepsilon\in(0,\pi]$ there is $N$ with $\sigma_i(H_n)>\pi-\varepsilon$ for all $n\ge N$, hence
$L_i\ge\pi-\varepsilon$. Letting $\varepsilon\downarrow0$ gives $L_i\ge\pi$. Therefore $L_i=\pi$:
\[
\lim_{n\to\infty}\sigma_i(H_n)=\pi\qquad\text{for every fixed }i\ge1 .
\]
The strict inequality $\sigma_i(H_n)<\pi$ for all finite $n$ (Lemma~\ref{lem:upper}) shows the
limit is approached but never attained.

\emph{Rate.} We prove the $O(1/\log n)$ upper estimate on $\pi-\sigma_i(H_n)$ directly from the
window construction. Fix $i$ and let $n$ be large. Choose the window ratio
$Q:=n^{1/(2i)}\ge2$ (valid for $n\ge 2^{2i}$). Running the construction of
Lemma~\ref{lem:lower} with this $Q$, the $i$ windows fit inside $\{1,\dots,N\}$ with
$N\le Q^{2i}=n$, so $U\subseteq V_n$, and
\[
\sigma_i(H_n)\ \ge\ \lambda_{\min}(B)\ \ge\ \Bigl(\pi-\frac{\pi+8}{\log Q}\Bigr)-iQ^{-1/2}
=\pi-\frac{2i(\pi+8)}{\log n}-i\,n^{-1/(4i)} .
\]
The last term decays super-polynomially compared to $1/\log n$, so for each fixed $i$,
\[
0<\pi-\sigma_i(H_n)\ \le\ \frac{2i(\pi+8)}{\log n}+i\,n^{-1/(4i)}=O\!\Bigl(\frac1{\log n}\Bigr),
\]
the lower bound $0<$ being Lemma~\ref{lem:upper}. This proves the rate stated in
Theorem~\ref{thm:main}.
\end{proof}

\section{Sharp rate for the largest eigenvalue}

\begin{proposition}\label{prop:rate}
For each fixed $i\ge1$, $\;0<\pi-\sigma_i(H_n)=O\!\left(\dfrac1{\log n}\right)$ as $n\to\infty$
(proved above). For the largest eigenvalue the sharp leading asymptotics are the de Bruijn--Wilf
relation
\[
\pi-\sigma_1(H_n)\ \sim\ \frac{\pi^3}{\log n}\qquad(n\to\infty),
\]
and consequently, for every fixed $i\ge1$,
$\;\pi-\sigma_i(H_n)=\Omega(1/\log n)$, so that $\pi-\sigma_i(H_n)=\Theta(1/\log n)$.
\end{proposition}

\begin{proof}
The upper bound $\pi-\sigma_i(H_n)=O(1/\log n)$ is the rate established in the proof of
Theorem~\ref{thm:main}. The sharp asymptotic
$\pi-\sigma_1(H_n)\sim\pi^3/\log n$ is the theorem of N.~G.~de Bruijn and H.~S.~Wilf,
\emph{On Hilbert's inequality in $n$ dimensions}, Bull.\ Amer.\ Math.\ Soc.\ \textbf{68} (1962),
70--73. For the lower bound at general $i$, note $\sigma_i(H_n)\le\sigma_1(H_n)$, hence
$\pi-\sigma_i(H_n)\ge\pi-\sigma_1(H_n)\sim\pi^3/\log n$, which gives
$\pi-\sigma_i(H_n)=\Omega(1/\log n)$. Combined with the $O(1/\log n)$ upper bound this yields
$\pi-\sigma_i(H_n)=\Theta(1/\log n)$ for every fixed $i$. The constant $\pi^3$ is asserted only for
$i=1$; we do not claim a matching constant for $i\ge2$.
\end{proof}

\begin{remark}
Finer corrections to $\pi-\sigma_1(H_n)$ (lower-order terms involving $\log\log n$) appear in the
subsequent literature. Direct numerical evaluation of $\sigma_1(H_n)$ for $n$ up to several
thousand gives values of $(\pi-\sigma_1(H_n))\log n$ around $4$--$5$, still well below
$\pi^3\approx31.0$; this reflects the large $\log\log n/\log n$ correction at such $n$ and is
consistent with the de Bruijn--Wilf expansion. Theorem~\ref{thm:main} does not rely on
Proposition~\ref{prop:rate}: its qualitative limit and its $O(1/\log n)$ rate both follow from the
self-contained Lemmas~\ref{lem:bounded}--\ref{lem:lower}.
\end{remark}

\end{document}
